\documentclass[11pt]{article}
\usepackage{mathblatt}


\blattfuss{Koerper}{Lösungen}
\begin{document}
\noindent{\large\bfseries Koerper \textperiodcentered{} Lösungen}\par\medskip
\begin{pfloesung}{}
\lz{1.}{$21$ cm²}{}{}
\end{pfloesung}
\begin{pfloesung}{}
\lz{2.}{$4\,000$ ml}{}{}
\end{pfloesung}
\begin{pfloesung}{}
\lz{3.}{$P : q$}{$r$}{}
\end{pfloesung}
\begin{pfloesung}{}
\lz{4.}{$9$}{}{}
\end{pfloesung}
\begin{pfloesung}{}
\lz{5.}{$48$ cm³}{$8 \cdot 6$}{}
\end{pfloesung}
\begin{pfloesung}{}
\lz{6.}{$9$ cm}{$18 : 2$}{}
\end{pfloesung}
\begin{pfloesung}{}
\lz{7.}{$3$ l}{}{}
\end{pfloesung}
\begin{pfloesung}{}
\lz{8.}{$4a^2$}{}{}
\end{pfloesung}
\begin{pfloesung}{}
\lz{9.}{Volumen}{}{}
\end{pfloesung}
\begin{pfloesung}{}
\lz{10.}{auch wenn das Prisma liegt}{Grundfläche: eines der beiden gleichen Dreiecke; die Deckfläche ist das andere}{}
\end{pfloesung}
\begin{pfloesung}{}
\lz{11.}{$36$ cm³}{$6 \cdot 3 \cdot 2$}{}
\end{pfloesung}
\begin{pfloesung}{}
\lz{12.}{$126$ cm³}{$18 \cdot 7$}{}
\end{pfloesung}
\begin{pfloesung}{}
\lz{13.}{$72$ cm³}{$V = 27 \cdot 8 : 3$}{}
\end{pfloesung}
\begin{pfloesung}{}
\lz{14.}{$125$ cm³}{$5 \cdot 5 \cdot 5$}{}
\end{pfloesung}
\begin{pfloesung}{}
\lz{15.}{$\tfrac{1}{3} \cdot \pi \cdot 3^2 \cdot 7 \approx 66{,}0$ cm³}{$V$}{}
\end{pfloesung}
\begin{pfloesung}{}
\lz{16.}{$\tfrac{4}{3} \cdot \pi \cdot 3^3 \approx 113{,}1$ cm³}{$V$}{}
\end{pfloesung}
\begin{pfloesung}{}
\lz{17.}{$\tfrac{4}{3} \cdot \pi \cdot 8^3 \approx 2\,144{,}7$ cm³}{$r = 8$ cm; $V$}{}
\end{pfloesung}
\begin{pfloesung}{}
\lz{18.}{$\tfrac{4}{3} \cdot \pi \cdot 2{,}5^3 \approx 65{,}4$ cm³}{$V$}{}
\end{pfloesung}
\begin{pfloesung}{}
\lz{19.}{$\tfrac{9 \pi}{2} \approx 14{,}1$ l}{$V = \tfrac{4}{3} \cdot \pi \cdot 15^3 \approx 14\,137{,}2$ cm³ }{}
\end{pfloesung}
\begin{pfloesung}{}
\lz{20.}{$105$ cm³}{Grundfläche $7 \cdot 3 = 21$ cm²; $21 \cdot 5$}{}
\end{pfloesung}
\begin{pfloesung}{}
\lz{21.}{$180$ cm³}{Grundfläche $8 \cdot 5 : 2 = 20$ cm²; $20 \cdot 9$}{}
\end{pfloesung}
\begin{pfloesung}{}
\lz{22.}{$120$ cm³}{Grundfläche $(7 + 3) : 2 \cdot 4 = 20$ cm²; $20 \cdot 6$}{}
\end{pfloesung}
\begin{pfloesung}{}
\lz{23.}{$360$ cm³}{$G = 6 \cdot 6 + 6 \cdot 3 : 2 = 45$ cm²; $V = 45 \cdot 8$}{}
\end{pfloesung}
\begin{pfloesung}{}
\lz{24.}{$63\,000$ cm³}{Grundfläche $90 \cdot 20 : 2 = 900$ cm²; $900 \cdot 70$}{}
\end{pfloesung}
\begin{pfloesung}{}
\lz{25.}{$84$ cm³}{$G = 6 \cdot 6 = 36$ cm²; $V = 36 \cdot 7 : 3$}{}
\end{pfloesung}
\begin{pfloesung}{}
\lz{26.}{$120$ cm³}{$G = 8 \cdot 5 = 40$ cm²; $V = 40 \cdot 9 : 3$}{}
\end{pfloesung}
\begin{pfloesung}{}
\lz{27.}{$200 \pi \approx 628{,}3$ cm³}{$\pi \cdot 5^2 \cdot 8$}{}
\end{pfloesung}
\begin{pfloesung}{}
\lz{28.}{das sind etwa $350$ ml}{$\pi \cdot 3{,}5^2 \cdot 9 \approx 346{,}4$ cm³ $\approx 346$ ml}{}
\end{pfloesung}
\begin{pfloesung}{}
\lz{29.}{$80 \pi \approx 251{,}3$ cm³}{$r = 4$ cm; $\pi \cdot 4^2 \cdot 5$}{}
\end{pfloesung}
\begin{pfloesung}{}
\lz{30.}{$\approx 147{,}3$ cm³}{$\pi \cdot 2{,}5^2 \cdot 7{,}5$}{}
\end{pfloesung}
\begin{pfloesung}{}
\lz{31.}{$\approx 5{,}3$ l}{$\pi \cdot 11^2 \cdot 14 \approx 5\,321{,}9$ cm³ }{}
\end{pfloesung}
\begin{pfloesung}{}
\lz{32.}{$\tfrac{3971}{150} \approx 26{,}5$ cm³}{$G = 3{,}8^2 = 14{,}44$ cm²; $V = 14{,}44 \cdot 5{,}5 : 3$}{}
\end{pfloesung}
\begin{pfloesung}{}
\lz{33.}{$\tfrac{1}{3} \cdot \pi \cdot 4^2 \cdot 9 \approx 150{,}8$ cm³}{$r = 4$ cm; $V$}{}
\end{pfloesung}
\begin{pfloesung}{}
\lz{34.}{$\tfrac{1}{3} \cdot \pi \cdot 2{,}5^2 \cdot 6{,}4 \approx 41{,}9$ cm³}{$V$}{}
\end{pfloesung}
\begin{pfloesung}{}
\lz{35.}{$\approx 415$ ml}{$V = \tfrac{1}{3} \cdot \pi \cdot 6^2 \cdot 11 \approx 414{,}7$ cm³ }{}
\end{pfloesung}
\begin{pfloesung}{{\small\color{mbgrau}P10 2026 FOR · 1f}}
\lz{36.}{27 cm³}{V = a³ = (3 cm)³ = 27 cm³}{}
\end{pfloesung}
\begin{pfloesung}{{\small\color{mbgrau}P10 2016 · 3c}}
\lz{37.}{288$\pi$ $\approx$ 904,8 cm³}{r = 6 cm; V = $\tfrac{4}{3}$ · $\pi$ · 6³}{}
\end{pfloesung}
\begin{pfloesung}{{\small\color{mbgrau}P10 2014 · 5a}}
\lz{38.}{862 125 cm³ ($\approx$ 0,86 m³)}{V = G · h = 5 225 cm² · 165 cm}{}
\end{pfloesung}
\begin{pfloesung}{{\small\color{mbgrau}P10 2022 · 2b}}
\lz{39.}{wahr, V = $\pi$ · 3,2² · 7 = 71,68$\pi$ $\approx$ 225,2 cm³ > 200 ml}{Grundfläche 10,24$\pi$ $\approx$ 32,2 cm²; 1 cm³ = 1 ml}{}
\end{pfloesung}
\begin{pfloesung}{{\small\color{mbgrau}P10 2021 · 4a}}
\lz{40.}{V = $\pi$ · 29² · 95 $\approx$ 250 998 cm³ $\approx$ 251 l}{1 l = 1 dm³ = 1000 cm³}{}
\end{pfloesung}
\begin{pfloesung}{{\small\color{mbgrau}P10 2023 · 5b}}
\lz{41.}{wahr, V = $\pi$ · 4² · 8 = 128$\pi$ $\approx$ 402,1 cm³ $\approx$ 400 ml}{Grundfläche 16$\pi$ $\approx$ 50,3 cm²; 1 cm³ = 1 ml}{}
\end{pfloesung}
\begin{pfloesung}{}
\lz{42.}{Radius $4$ cm}{Durchmesser}{}
\end{pfloesung}
\begin{pfloesung}{}
\lz{43.}{Prisma}{}{}
\end{pfloesung}
\begin{pfloesung}{}
\lz{44.}{$6$}{$X = 5$, $Y = 4$, $Z$}{}
\end{pfloesung}
\begin{pfloesung}{}
\lz{45.}{$8$ cm}{Kreis: $r = 3$ cm; Kreisausschnitt: Radius ; die Mantellinie}{}
\end{pfloesung}
\begin{pfloesung}{}
\lz{46.}{Quader}{}{}
\end{pfloesung}
\begin{pfloesung}{}
\lz{47.}{zwischen B und D liegt genau eine Fläche}{D}{}
\end{pfloesung}
\begin{pfloesung}{}
\lz{48.}{ja}{}{}
\end{pfloesung}
\begin{pfloesung}{}
\lz{49.}{ein Quadrat als Grundfläche, vier Dreiecke laufen in einer Spitze zusammen}{Pyramide}{}
\end{pfloesung}
\begin{pfloesung}{}
\lz{50.}{$4 \pi \approx 12{,}6$ cm, also $5$ cm × $12{,}6$ cm}{Umfang $2 \cdot \pi \cdot 2$}{}
\end{pfloesung}
\begin{pfloesung}{{\small\color{mbgrau}P10 2017 · 1c}}
\lz{51.}{Prisma}{}{}
\end{pfloesung}
\begin{pfloesung}{{\small\color{mbgrau}P10 2018 · 1i}}
\lz{52.}{Pyramide}{}{}
\end{pfloesung}
\begin{pfloesung}{{\small\color{mbgrau}P10 2016 · 1j}}
\lz{53.}{mittleres Quadrat der senkrechten Dreierreihe}{}{}
\end{pfloesung}
\begin{pfloesung}{{\small\color{mbgrau}P10 2022 · 2a}}
\lz{54.}{drittes Netz; Länge 6,4$\pi$ $\approx$ 20,1 cm, Breite 7 cm}{Länge = Umfang 2 · $\pi$ · 3,2; Breite = h = 7 cm}{}
\end{pfloesung}
\begin{pfloesung}{}
\lz{55.}{Mantellinie}{}{}
\end{pfloesung}
\begin{pfloesung}{}
\lz{56.}{Seitenhöhe}{}{}
\end{pfloesung}
\begin{pfloesung}{}
\lz{57.}{$6$ cm}{$90 : (5 \cdot 3)$}{}
\end{pfloesung}
\begin{pfloesung}{}
\lz{58.}{$27\,000$}{$30$ cm, denn $30 \cdot 30 \cdot 30$}{}
\end{pfloesung}
\begin{pfloesung}{}
\lz{59.}{$6$ cm}{$210 : 35$}{}
\end{pfloesung}
\begin{pfloesung}{}
\lz{60.}{$8$ cm}{$h = 3 \cdot 96 : 36$}{}
\end{pfloesung}
\begin{pfloesung}{}
\lz{61.}{$3\,000 : (\pi \cdot 10^2) \approx 9{,}5$ cm}{$r = 10$ cm, $V = 3\,000$ cm³; $h$}{}
\end{pfloesung}
\begin{pfloesung}{}
\lz{62.}{$\approx 6{,}5$ cm}{$\sqrt{2\,000 : (\pi \cdot 15)}$; erst teilen, dann die Wurzel}{}
\end{pfloesung}
\begin{pfloesung}{}
\lz{63.}{$\sqrt{3 \cdot 300 : (\pi \cdot 9)} \approx 5{,}6$ cm}{$r$}{}
\end{pfloesung}
\begin{pfloesung}{}
\lz{64.}{$\sqrt[3]{3 \cdot 1\,000 : (4 \cdot \pi)} \approx 6{,}2$ cm}{$r$}{}
\end{pfloesung}
\begin{pfloesung}{}
\lz{65.}{$3 \cdot 150 : (\pi \cdot 4^2) \approx 9{,}0$ cm}{$h$}{}
\end{pfloesung}
\begin{pfloesung}{{\small\color{mbgrau}P10 2023 · 5d}}
\lz{66.}{h = 125/(2$\pi$) $\approx$ 19,9 cm}{1 l = 1000 cm³; h = V : ($\pi$ · r²) = 1000 : (16$\pi$); Grundfläche 16$\pi$ $\approx$ 50,3 cm²}{}
\end{pfloesung}
\begin{pfloesung}{{\small\color{mbgrau}P10 2022 · 2d}}
\lz{67.}{r = √(425/(7$\pi$)) $\approx$ 4,40 cm}{r² = 425 : ($\pi$ · 7) = 425/(7$\pi$) $\approx$ 19,33}{}
\end{pfloesung}
\begin{pfloesung}{}
\lz{68.}{Durchmesser}{}{}
\end{pfloesung}
\begin{pfloesung}{}
\lz{69.}{Quader und Dreiecksprisma}{Wohnteil; Dach}{}
\end{pfloesung}
\begin{pfloesung}{}
\lz{70.}{$96$ dm³}{$8 \cdot 4 \cdot 2 + 8 \cdot 2 \cdot 2 = 64 + 32$}{}
\end{pfloesung}
\begin{pfloesung}{}
\lz{71.}{$\approx 1{,}76$ kg}{außen $r = 1{,}1$ cm, innen $r = 1{,}0$ cm; $V = \pi \cdot (1{,}1^2 - 1^2) \cdot 300 \approx 197{,}9$ cm³; $m \approx 1\,761$ g }{}
\end{pfloesung}
\begin{pfloesung}{}
\lz{72.}{$a : 2$, die Kugel hat $\tfrac{\pi}{6} \cdot a^3$}{a) $r = 6$ cm; $V \approx 904{,}8$ cm³. b) Würfel $1\,728$ cm³; Abfall $\approx 823{,}2$ cm³ $\approx 47{,}6$ \%. c) Mit Kante $a$ ist $r$}{}
\end{pfloesung}
\begin{pfloesung}{}
\lz{73.}{$2\,142{,}9$ cm³}{Karton $24 \cdot 24 \cdot 9 = 5\,184$ cm³; Torte $\pi \cdot 11^2 \cdot 8 \approx 3\,041{,}1$ cm³; $5\,184 - 3\,041{,}1$}{}
\end{pfloesung}
\begin{pfloesung}{}
\lz{74.}{$\approx 26\,832$ cm³}{Zylinder $\pi \cdot 18{,}5^2 \cdot 36 \approx 38\,708$ cm³; Kegel $\tfrac{1}{3} \cdot \pi \cdot 18^2 \cdot 35 \approx 11\,875$ cm³; Rest }{}
\end{pfloesung}
\begin{pfloesung}{{\small\color{mbgrau}P10 2024 · 4c}}
\lz{75.}{51 350,25$\pi$ $\approx$ 161 322 cm³}{V\_Zylinder = $\pi$ · 30,5² · 81 = 75 350,25$\pi$ $\approx$ 236 719,8 cm³; V\_Kegel = $\tfrac{1}{3}$ · $\pi$ · 30² · 80 = 24 000$\pi$ $\approx$ 75 398,2 cm³}{}
\end{pfloesung}
\begin{pfloesung}{}
\lz{76.}{$50$ cm und $70$ cm}{}{}
\end{pfloesung}
\begin{pfloesung}{}
\lz{77.}{Höhe von der Spitze senkrecht zur Mitte der Grundfläche, von dort zur Mitte einer Grundkante , zurück zur Spitze}{halbe Grundkante, $3$ m}{}
\end{pfloesung}
\begin{pfloesung}{}
\lz{78.}{$108$ cm²}{$M = 4 \cdot 6 \cdot 9 : 2$}{}
\end{pfloesung}
\begin{pfloesung}{}
\lz{79.}{$60$ m²}{$V = 6 \cdot 6 \cdot 4 : 3 = 48$ m³; Glas: $M = 4 \cdot 6 \cdot 5 : 2$}{}
\end{pfloesung}
\begin{pfloesung}{}
\lz{80.}{$72 \pi \approx 226{,}2$ cm²}{$2 \cdot \pi \cdot 4 \cdot 9$}{}
\end{pfloesung}
\begin{pfloesung}{}
\lz{81.}{$140 \pi \approx 439{,}8$ cm²}{$2 \cdot \pi \cdot 5^2 + 2 \cdot \pi \cdot 5 \cdot 9$}{}
\end{pfloesung}
\begin{pfloesung}{}
\lz{82.}{$\pi \cdot 4 \cdot 11 \approx 138{,}2$ cm²}{$M$}{}
\end{pfloesung}
\begin{pfloesung}{}
\lz{83.}{$30 \pi \approx 94{,}2$ cm²}{$G = \pi \cdot 3^2 \approx 28{,}27$ cm², $M = \pi \cdot 3 \cdot 7 \approx 65{,}97$ cm², $O$}{}
\end{pfloesung}
\begin{pfloesung}{}
\lz{84.}{$\approx 23{,}1$ m}{$M = \pi \cdot 3{,}5 \cdot 6{,}2 \approx 68{,}2$ m²; Kosten $\approx 3\,068$ €; Dachhöhe $= \sqrt{6{,}2^2 - 3{,}5^2} \approx 5{,}12$ m; Gesamthöhe }{}
\end{pfloesung}
\begin{pfloesung}{{\small\color{mbgrau}P10 2018 · 6b}}
\lz{85.}{138,88$\pi$ $\approx$ 436,3 m²}{M = $\pi$ · d · h = $\pi$ · 6,4 · 21,7 = 138,88$\pi$ m²}{}
\end{pfloesung}
\begin{pfloesung}{}
\lz{86.}{$3$ cm}{Durchmesser; $r$}{}
\end{pfloesung}
\begin{pfloesung}{}
\lz{87.}{$28$ cm}{Radius; $d$}{}
\end{pfloesung}
\begin{pfloesung}{}
\lz{88.}{vierfach, weil der Radius im Quadrat steht}{A: $\pi \cdot 3^2 \cdot 5 = \pi \cdot 45$; B: $\pi \cdot 6^2 \cdot 5 = \pi \cdot 180$; $180 : 45 = 4$}{}
\end{pfloesung}
\begin{pfloesung}{}
\lz{89.}{$\tfrac{3}{4}$}{alt: $\tfrac{1}{3} \cdot 36 \cdot 5 = 60$ cm³; neu: $\tfrac{1}{3} \cdot 9 \cdot 15 = 45$ cm³; kleiner, $\tfrac{45}{60}$; Grundfläche ein Viertel, Höhe dreifach}{}
\end{pfloesung}
\begin{pfloesung}{}
\lz{90.}{das Achtfache, weil $2^3 = 8$}{A: $V \approx 33{,}5$ cm³, B: $V \approx 268{,}1$ cm³; $268{,}1 : 33{,}5 \approx 8$}{}
\end{pfloesung}
\begin{pfloesung}{}
\lz{91.}{$15 \pi \approx 47{,}1$ cm³}{Zylinder: $\pi \cdot 3^2 \cdot 5 \approx 141{,}4$ cm³; Kegel: $\tfrac{1}{3}$ davon }{}
\end{pfloesung}
\begin{pfloesung}{}
\lz{92.}{Der Radius steht im Quadrat: Der doppelte Radius gibt eine viermal so große Grundfläche, die Höhe bleibt gleich}{}{}
\end{pfloesung}
\begin{pfloesung}{}
\lz{93.}{Das Kegelvolumen ist ein Drittel von Grundfläche mal Höhe, das Zylindervolumen Grundfläche mal Höhe}{}{}
\end{pfloesung}
\begin{pfloesung}{}
\lz{94.}{$12 \cdot \pi \cdot r^2 \cdot h$, also das Zwölffache des Zylindervolumens $\pi \cdot r^2 \cdot h$}{Die Aussage trifft nicht zu; Mit dem Radius $r$ des Zylinders und der Höhe $h$: $V_K = \tfrac{1}{3} \cdot \pi \cdot (6r)^2 \cdot h$}{}
\end{pfloesung}
\begin{pfloesung}{{\small\color{mbgrau}P10 2021 · 4b}}
\lz{95.}{nein, Volumen vervierfacht sich (319,58$\pi$ l $\approx$ 1004 l)}{(2r)² = 4r²}{}
\end{pfloesung}
\begin{pfloesung}{}
\lz{96.}{sechs Rechtecke, je zwei gleich: $4 \times 2$, $4 \times 1$ und $2 \times 1$ , so aneinander, dass sie zu einem Quader falten}{in cm}{}
\end{pfloesung}
\begin{pfloesung}{}
\lz{97.}{Es fehlen die zwei Dreiecke mit den Seiten $4$ cm, $5$ cm und $6$ cm, je eins oben und unten an einem Rechteck}{}{}
\end{pfloesung}
\begin{pfloesung}{}
\lz{98.}{zwei Dreiecke mit den Seiten $8$ cm, $5$ cm, $5$ cm}{}{}
\end{pfloesung}
\begin{pfloesung}{}
\lz{99.}{alle vier Seiten des Quadrats $4$ cm}{}{}
\end{pfloesung}
\begin{pfloesung}{}
\lz{100.}{$20 \pi \approx 62{,}8$ cm³}{$\pi \cdot 2^2 \cdot 5$}{}
\end{pfloesung}
\begin{pfloesung}{{\small\color{mbgrau}P10 2020 · 5a}}
\lz{101.}{Rückfläche 15 cm × 12 cm rechts an die Grundfläche; zweites Trapez unten an die Grundfläche}{}{}
\end{pfloesung}
\begin{pfloesung}{{\small\color{mbgrau}P10 2019 · 4a}}
\lz{102.}{Netz mit drei Rechtecken und zwei gleichseitigen Dreiecken; a = 1,50 m}{}{}
\end{pfloesung}
\begin{pfloesung}{{\small\color{mbgrau}P10 2014 · 5b}}
\lz{103.}{Netz um Rechtecke 57 cm × 165 cm und 110 cm × 165 cm ergänzt; a = 110 cm, b = 57 cm, d = 165 cm}{}{}
\end{pfloesung}
\begin{pfloesung}{{\small\color{mbgrau}P10 2023 · 5a}}
\lz{104.}{Rechteck, a = 8$\pi$ $\approx$ 25,1 cm, b = 8 cm}{a = 2 · $\pi$ · r (Umfang); b = h}{}
\end{pfloesung}
\begin{pfloesung}{}
\lz{105.}{Tiefe $4$ cm schräg als $2$ cm}{Vorderseite $6$ cm $\times$ $3$ cm in wahrer Größe; Kästchendiagonalen}{}
\end{pfloesung}
\begin{pfloesung}{}
\lz{106.}{nicht eine Seitenkante und nicht die Höhe einer Seitenfläche}{Diagonalen der Grundfläche zeichnen, ihren Schnittpunkt gestrichelt mit der Spitze verbinden: das ist die Höhe, $7$ cm}{}
\end{pfloesung}
\begin{pfloesung}{}
\lz{107.}{Spitze in der Mitte des Deckels}{Grundkreis auf dem Boden, er berührt alle vier Seitenwände , im Schrägbild als Ellipse; Durchmesser $50$ cm $=$ Kartonbreite; Schnitt der Diagonalen}{}
\end{pfloesung}
\begin{pfloesung}{}
\lz{108.}{Quadrat mit $5$ cm Seitenlänge, an jeder Seite ein gleichschenkliges Dreieck mit Grundseite $5$ cm und Höhe $7$ cm}{}{}
\end{pfloesung}
\begin{pfloesung}{}
\lz{109.}{Höhe $5$ cm senkrecht vom Mittelpunkt, gestrichelt}{Grundkreis als flache Ellipse, $4$ cm breit, mit Mittelpunkt}{}
\end{pfloesung}
\begin{pfloesung}{}
\lz{110.}{oben der Kreis als Ellipse mit dem Radius $2{,}5$ cm}{Kegel mit der Spitze unten}{}
\end{pfloesung}
\begin{pfloesung}{}
\lz{111.}{flache Ellipse durch den Mittelpunkt als Äquator}{Kreis mit dem Radius $2$ cm; hintere Hälfte gestrichelt}{}
\end{pfloesung}
\begin{pfloesung}{}
\lz{112.}{Höhe $5$ cm senkrecht darüber, gestrichelt}{Grundquadrat: vordere Kante $4$ cm, Tiefenkanten $2$ cm unter 45°; Diagonalen schneiden sich im Fußpunkt}{}
\end{pfloesung}
\begin{pfloesung}{{\small\color{mbgrau}P10 2024 · 4b}}
\lz{113.}{Skizze: Ellipse auf der Bodenfläche, Spitze mittig oben; 61 cm × 61 cm × 81 cm; 31 cm schmaler als Durchmesser 60 cm, 101 cm Höhe unnötig groß}{Durchmesser 2 · 30 = 60 cm; Höhe 80 cm}{}
\end{pfloesung}
\begin{pfloesung}{{\small\color{mbgrau}P10 2015 · 6b}}
\lz{114.}{Höhe von der Spitze zum Diagonalenschnittpunkt, h = 0,68 m; Grundkante 0,80 m}{}{}
\end{pfloesung}
\begin{pfloesung}{\textbf{Prüfstein} \textperiodcentered{} P10 2026 FOR · 2}
\lz{a)}{552,25$\pi$ $\approx$ 1734,9 m³}{V = $\pi$ · r² · h = $\pi$ · 4,7² · 25 = 552,25$\pi$ m³; r² = 22,09 m²}{}
\lz{b)}{$\approx$ 2539,66 € (rund 2540 €)}{M = $\pi$ · r · s = 40,42$\pi$ $\approx$ 126,98 m²}{}
\lz{c)}{25 + √51,87 $\approx$ 32,2 m}{Kegelhöhe = √(8,6² $-$ 4,7²) = √51,87 $\approx$ 7,20 m}{}
\lz{d)}{falsch, V\_Kegel = $\tfrac{1}{3}$ · $\pi$ · (3r)² · h = 3 · $\pi$ · r² · h = 3 · V\_Zylinder $\Rightarrow$ dreifaches Volumen}{(3r)² = 9r²; 9 : 3 = 3}{}
\end{pfloesung}
\end{document}